\section{Geometry near a node}\label{sec:node}

Throughout, $\dc=\tfrac i2(\bar\partial-\partial)$, so that
$d\dc=i\partial\bar\partial$, and for a K\"ahler form $\omega$ and a real
function $f$ the contraction $\Lambda_\omega d\dc f$ equals $-\tfrac12\Delta f$,
where $\Delta=d^*d$. The Hodge Laplacian on forms is $\Delta=dd^*+d^*d$.
The letter $C$ denotes constants independent of the family parameter $t$
and of the auxiliary radii $R$ and $A$ introduced below; $c(R,A)$ denotes
constants that may depend on them.

\subsection{Normal form}
Let $p$ be a singular point of $X_0$. Since $\cX$ is smooth and the fibre $\pi^{-1}(0)$ is reduced, $\pi$ is
a local equation of $X_0$ at $p$; as $X_0$ has an ordinary double point
there, $\pi$ has a nondegenerate critical point at $p$.

\begin{lemma}\label{lem:normal-form}
There are holomorphic coordinates $w=(w_1,\ldots,w_4)$ on a neighbourhood
$U_p$ of $p$ in $\cX$, centred at $p$, and numbers $\sigma_1,\ldots,\sigma_4>0$,
such that
\begin{equation}\label{eq:normal-form}
 \pi=Q_\sigma(w):=\sum_{j=1}^4\sigma_jw_j^2,\qquad
 \omega=\omega_{\rm fl}+\omega',\quad
 \omega_{\rm fl}=\tfrac i2\sum_jdw_j\wedge d\bar w_j,
\end{equation}
where $\omega'$ is a smooth $(1,1)$-form on $U_p$ vanishing at $p$.
\end{lemma}
\begin{proof}
By the holomorphic Morse lemma there are coordinates $z$ with
$\pi=\sum_jz_j^2$. Let $\omega_p=\tfrac i2\sum h_{j\bar k}dz_j\wedge d\bar z_k$
at $p$, with $(h_{j\bar k})$ positive Hermitian, and choose an invertible
matrix $T$ such that the linear change $z=Tw'$ brings $\omega_p$ to
$\tfrac i2\sum dw'_j\wedge d\bar w'_j$. In these coordinates
$\pi=(w')^{\mathsf t}Sw'$ with $S=T^{\mathsf t}T$ complex symmetric and invertible. By the
Autonne--Takagi factorization \cite[Section~4.4]{horn-johnson} there is
a unitary $U$ with $S=U^{\mathsf t}\operatorname{diag}(\sigma_j)U$ and
$\sigma_j\ge0$; invertibility gives $\sigma_j>0$. Put $w=Uw'$.
\end{proof}

We fix such coordinates at every node $p_i$, $1\le i\le N$, with coefficients
$\sigma^{(i)}$, and write $r=|w|$ in each chart. Thus, in the chart,
\begin{equation}\label{eq:local-fibres}
 X_t\cap U_p=\{Q_\sigma=t\}\cap U_p\qquad(t\in\DD),
\end{equation}
and $X_0\cap U_p$ is an open subset of the cone
$C_\sigma=\{Q_\sigma=0\}\subset\CC^4$.

\subsection{The cone and its link}
Let $g_C$ be the metric induced on $C_\sigma\setminus\{0\}$ by the flat
metric of $\CC^4$. It is a K\"ahler cone metric $g_C=dr^2+r^2h$ with
K\"ahler form $\omega_C=\omega_{\rm fl}|_{C_\sigma}$ and link
$L_\sigma=C_\sigma\cap S^7$ with the induced metric $h$. The circle
$\{e^{i\phi}\}$ acts freely and isometrically on $L_\sigma$ with orbits of
length $2\pi$, and the quotient is the smooth quadric surface
$Z_\sigma=\{Q_\sigma=0\}\subset\PP^3$. Let $\varpi:L_\sigma\to Z_\sigma$ be the
projection. It is a Riemannian submersion onto $Z_\sigma$ with the
K\"ahler form $\omega_Z$ given by $\varpi^*\omega_Z=\tfrac i2\partial\bar\partial
\log|w|^2$, and the one-form $\xi=\dc\log|w|^2$ restricts to $d\phi$ on
the orbits and satisfies $d\xi=2\varpi^*\omega_Z$.

The surface $Z_\sigma$ is isomorphic to $\PP^1\times\PP^1$. Let $F_1,F_2\in
H^2(Z_\sigma;\ZZ)$ be the classes of the fibres of its two rulings, so that
$F_1^2=F_2^2=0$, $F_1F_2=1$, and the hyperplane class is $F_1+F_2$. The link is the
unit circle bundle of $\OO(-1)|_{Z_\sigma}$, and the Gysin sequence gives
\begin{equation}\label{eq:link-cohomology}
 H^1(L_\sigma;\ZZ)=0,\qquad
 H^2(L_\sigma;\ZZ)=H^2(Z_\sigma;\ZZ)/\ZZ(F_1+F_2)\cong\ZZ,\qquad
 H^3(L_\sigma;\ZZ)\cong\ZZ .
\end{equation}
The group $H^2(L_\sigma;\ZZ)$ is generated by $\varpi^*F_1=-\varpi^*F_2$. Which of the
two generators is used at a node is fixed in Section~\ref{sec:harmonic}
by the small resolution.

\begin{lemma}\label{lem:link-form}
There is a unique harmonic two-form $\vartheta$ on $(L_\sigma,h)$ in the class $\varpi^*F_1$. It satisfies
\begin{equation}\label{eq:link-constant}
 \|\vartheta\|^2_{L^2(L_\sigma,h)}=\pi
\end{equation}
for every $\sigma$. Its pull-back to $C_\sigma\setminus\{0\}$ under the
radial projection, again denoted $\vartheta$, is a closed and coclosed
real primitive $(1,1)$-form for $(g_C,\omega_C)$, with
$|\nabla_C^j\vartheta|_{g_C}=O(r^{-2-j})$ for all $j\ge0$ and vanishing
contraction with $\partial_r$. The form $*_h\vartheta$ is closed and
generates $H^3(L_\sigma;\RR)$.
\end{lemma}
\begin{proof}
Let $\vartheta_Z$ be the harmonic representative on $(Z_\sigma,\omega_Z)$
of the class $e=\tfrac12(F_1-F_2)$. Since $H^{2,0}(Z_\sigma)=0$ it is a real
$(1,1)$-form, and it is primitive because $e\cdot(F_1+F_2)=0$ and $\omega_Z$
is a multiple of the hyperplane class. On a K\"ahler surface a primitive
real $(1,1)$-form is anti-self-dual, so
\[
 \int_{Z_\sigma}|\vartheta_Z|^2\,dV=-\int_{Z_\sigma}\vartheta_Z\wedge
 \vartheta_Z=-e^2=\tfrac12 .
\]
Put $\vartheta=\varpi^*\vartheta_Z$. It is closed, and its class is
$\varpi^*e=\varpi^*F_1$ because $\varpi^*(F_1+F_2)=0$. For the Riemannian submersion $\varpi$ with
unit vertical one-form $\xi$, the Hodge star of the pull-back of a
two-form $\beta$ is $\varpi^*(*_Z\beta)\wedge\xi$ for the orientation
$\varpi^*dV_Z\wedge\xi$. Hence $*_h\vartheta=-\vartheta\wedge\xi$, and
$d(\xi\wedge\vartheta)=2\varpi^*(\omega_Z\wedge\vartheta_Z)=0$ by
primitivity. Thus $\vartheta$ is harmonic; uniqueness follows from
$b_2(L_\sigma)=1$. Since the fibres of $\varpi$ have length $2\pi$,
\[
 \|\vartheta\|^2_{L^2(L_\sigma)}=2\pi\int_{Z_\sigma}|\vartheta_Z|^2\,dV=\pi .
\]
On the cone, $\vartheta$ is the pull-back of $\vartheta_Z$ under the
holomorphic map $C_\sigma\setminus\{0\}\to Z_\sigma$, hence closed, real,
of type $(1,1)$ and homogeneous of degree zero, which gives the stated
bounds and $\iota_{\partial_r}\vartheta=0$. With
$\omega_C=\tfrac12d(r^2\xi)=r\,dr\wedge\xi+r^2\varpi^*\omega_Z$ one finds
$\vartheta\wedge\omega_C^2=0$, because
$\vartheta_Z\wedge\omega_Z=0$ and $\vartheta\wedge \varpi^*\omega_Z^2$ is the
pull-back of a six-form from a surface. So $\vartheta$ is primitive, and
then $*\vartheta=-\vartheta\wedge\omega_C$ is closed. Finally
$\int_{L_\sigma}\vartheta\wedge*_h\vartheta=\pi\ne0$, so the closed form
$*_h\vartheta$ is not exact.
\end{proof}

\subsection{The smoothing cap}
Let $V_\sigma=\{Q_\sigma=1\}\subset\CC^4$ with the metric $g_V$ induced by
the flat metric. It is diffeomorphic to $T^*S^3$; in particular
\begin{equation}\label{eq:cap-topology}
 H^1(V_\sigma)=H^2(V_\sigma)=0,\qquad
 H^3(V_\sigma;\RR)\xrightarrow{\ \sim\ }H^3(L_\sigma;\RR),
\end{equation}
the second map being restriction to a large sphere. For $t=\epsilon^2u^2$
with $\epsilon>0$ and $|u|=1$,
\begin{equation}\label{eq:cap-scaling}
 \{Q_\sigma=t\}=\epsilon u\cdot V_\sigma ,
\end{equation}
and multiplication by $\epsilon u$ multiplies the flat metric by
$\epsilon^2$. Throughout we write
\begin{equation}\label{eq:scale}
 \epsilon=|t|^{1/2},\qquad u^2=t/|t| .
\end{equation}

\begin{lemma}\label{lem:cap-ac}
Put $n(w)=\overline{\partial Q_\sigma(w)}/|\partial Q_\sigma(w)|^2$, where
$\partial Q_\sigma=(2\sigma_jw_j)_j$. There are $A_0>0$ and a smooth
complex-valued function $\tau$ on $C_\sigma\cap\{r>A_0\}$ with
$\tau=1+O(r^{-2})$ such that
\[
 \Psi(w)=w+\tau(w)\,n(w)
\]
is a diffeomorphism from $C_\sigma\cap\{r>A_0\}$ onto the complement of a
compact subset of $V_\sigma$, and
\begin{equation}\label{eq:cap-ac}
 |\nabla_C^j(\Psi^*g_V-g_C)|_{g_C}\le C_jr^{-2-j}\qquad(j\ge0).
\end{equation}
\end{lemma}
\begin{proof}
Since $Q_\sigma$ is quadratic, $Q_\sigma(w+\tau n)=Q_\sigma(w)+\tau
+\tau^2Q_\sigma(n)$, because $dQ_\sigma(w)\cdot n(w)=1$. The vector
field $n$ is homogeneous of degree $-1$ under positive dilations and
$|n|\le Cr^{-1}$, so $Q_\sigma(n)=O(r^{-2})$ with corresponding bounds on
derivatives, and the equation $\tau+\tau^2Q_\sigma(n)=1$ has a
unique smooth solution near $1$ for large $r$, with
$\tau-1=O(r^{-2})$. The displacement $\tau n$ and its covariant
derivatives of order $j$ are $O(r^{-1-j})$, so $\Psi$ differs from the
inclusion of the cone by a map whose first derivatives are $O(r^{-2})$;
this gives \eqref{eq:cap-ac}. The same estimate shows that $\Psi$ is an
injective immersion for large $r$, and its image contains every point
$v\in V_\sigma$ of large norm: solving $Q_\sigma(v-\tau'n(v))=0$ in the
same way gives the cone point over $v$, up to a correction that is
removed by the inverse function theorem.
\end{proof}

\subsection{Conic surgery models}\label{sec:models}
Let $M_0=X_0\setminus\{p_1,\ldots,p_N\}$ with the K\"ahler metric $g_0$
induced by $\omega$. By \eqref{eq:normal-form}, on the end of $M_0$ at
$p_i$,
\begin{equation}\label{eq:central-rate}
 |\nabla_C^j(g_0-g_C)|_{g_C}\le C_jr^{1-j}\qquad(j\ge0),
\end{equation}
and likewise for the K\"ahler forms. Fix $R_0>0$ such that the balls
$\{r\le4R_0\}$ at the nodes lie in the charts and are pairwise disjoint.

Fix $0<R\le R_0$ and $A\ge A_0$. Let $g_0^R$ be a smooth metric on $M_0$
equal to $g_0$ where $r\ge2R$ and to $g_C$ where $r\le R$ at every node,
obtained by interpolation with a fixed cutoff in $r/R$. On each
$V_{\sigma^{(i)}}$ let $g_V^A$ be a smooth metric equal to $g_V$ on the
compact set bounded by $\Psi(\{r=A\})$ and to $(\Psi^{-1})^*g_C$ outside
$\Psi(\{r=2A\})$, interpolated in the same way. By \eqref{eq:central-rate}
and \eqref{eq:cap-ac},
\begin{equation}\label{eq:interpolation-errors}
 |g_0^R-g_0|_{g_0}\le CR,\qquad |g_V^A-g_V|_{g_V}\le CA^{-2}.
\end{equation}
These metrics need not be K\"ahler.

For $T>A_0$ let $V_\sigma^{T}$ be the open subset of $V_\sigma$ with
compact closure bounded by $\Psi(\{r=T\})$.

\begin{definition}\label{def:model}
Let $8A\epsilon<R$ and $|u|=1$. The \emph{conic surgery model}
$(\overline X_{\epsilon,u},\bar g_{\epsilon,u}^{R,A})$ is the closed
Riemannian manifold obtained from the central piece
$(M_0\cap\{r>2A\epsilon\},g_0^R)$ and, for each node, the cap core
$(V_{\sigma}^{4A},\epsilon^2g_V^A)$ with $\sigma=\sigma^{(i)}$, by
identifying each point $w$ of the central piece with
$2A\epsilon<|w|<4A\epsilon$ in the chart at the node with the point
$\Psi(w/\epsilon u)$ of the cap core.
\end{definition}

The identification is an isometry between open collars on which both
metrics are exactly conical, so the model is a smooth closed Riemannian
manifold. At each node the part of the model inside a section
$\{r=T\}$, $4A\epsilon<T<R$, is the union of the cap core and the
conical region $\{2A\epsilon<r\le T\}$; the map
$v\mapsto\epsilon u\,\Psi^{-1}(v)$ on the conical part of the cap
identifies it isometrically with the truncated cap
$(\overline{V_\sigma^{T/\epsilon}},\epsilon^2g_V^A)$.

\begin{lemma}\label{lem:condition-N}
Let $t=\epsilon^2u^2$. For $|t|$ small there are diffeomorphisms
$\Phi_t:\overline X_{\epsilon,u}\to X_t$ with
\begin{equation}\label{eq:condition-N}
 \limsup_{t\to0}\ \sup_{\overline X_{\epsilon,u}}
 \bigl|\Phi_t^*g_t-\bar g_{\epsilon,u}^{R,A}\bigr|_{\bar g_{\epsilon,u}^{R,A}}
 \le C\,(R+A^{-2}).
\end{equation}
Moreover, for unit complex numbers $u,u'$ there is a diffeomorphism
$\overline X_{\epsilon,u}\to\overline X_{\epsilon,u'}$ under which the two
model metrics differ by at most $C|u-u'|$ relative to either of them; here and
below, two metrics $g,\bar g$ differ by at most $\delta$ relative to
$\bar g$ if $(1-\delta)\bar g\le g\le(1+\delta)\bar g$.
\end{lemma}
\begin{proof}
On the cap core at a node define $\Phi_t(v)=\epsilon uv$, which lies in
$\{Q_\sigma=t\}=X_t\cap U_p$ by \eqref{eq:cap-scaling}. On the conical
region $\{2A\epsilon<|w|<R\}$ of the central piece put
\[
 \Psi_t(w)=\epsilon u\,\Psi(w/\epsilon u)=w+t\,\tau_t(w)\,n(w),
 \qquad\tau_t+t\tau_t^2Q_\sigma(n(w))=1 ,
\]
which is compatible with the gluing of the model, takes values in
$\{Q_\sigma=t\}$, and depends smoothly on $(t,w)$ for $|w|\ge R/2$ and
small $t$, with $\Psi_0$ the identity. The map $(t,w)\mapsto\Psi_t(w)$
is an open embedding of $\DD\times\{R/2<|w|<R\}$ into $\cX$ over
$\DD$, for a smaller disc. Push the two real coordinate vector fields of
$\DD$ forward by it, and extend them, by a partition of unity, to lifts
of these vector fields on a neighbourhood of $M_0\cap\{r\ge R/2\}$ in
$\cX$, where $\pi$ is a submersion. Flowing along the two lifts
successively extends $\Psi_t$ to embeddings of $M_0\cap\{r>R/2\}$ into
$X_t$, smooth in $t$ and equal to the identity at $t=0$, whose image is
the complement of the image of $\{|w|\le R/2\}$ and of the cap cores.
This defines $\Phi_t$.

On the exterior $\{r\ge R\}$ the metrics $\Phi_t^*g_t$ converge uniformly
to $g_0$, and $g_0^R$ differs from $g_0$ by at most $CR$ by
\eqref{eq:interpolation-errors}. On the cap core and the conical region, identified with a truncated
cap as after Definition~\ref{def:model}, the map $\Phi_t$ is
$v\mapsto\epsilon uv$; the form $\Phi_t^*\omega_{\rm fl}$
restricts to $\epsilon^2$ times the K\"ahler form of $g_V$, and by
\eqref{eq:normal-form} the form $\omega'$ is bounded by $C|w|\le CR$
relative to $\omega_{\rm fl}$ on the region $|w|\le R$ containing the
image. Hence $\Phi_t^*g_t$ differs from $\epsilon^2g_V$ by at most $CR$
there, and $g_V^A$ differs from $g_V$ by at most $CA^{-2}$. This proves
\eqref{eq:condition-N}.

For the last assertion write $u'=e^{i\delta}u$ and let $\chi$ be a fixed
cutoff, equal to one for $r\le R/2$ and to zero for $r\ge R$. The map of
the central piece equal to $w\mapsto e^{i\delta\chi(|w|)}w$ in the charts
and to the identity elsewhere preserves the cone $C_\sigma$, and together
with the identity on the cap cores it is compatible with the two gluing
maps, which involve only the region $\{r<4A\epsilon\}$, where $\chi=1$. Multiplication by a unit scalar is an
isometry of $g_C$, and the radial dependence of the angle changes the
metric $g_0^R$ by at most $C|\delta|$.
\end{proof}
